\documentclass[12pt,a4paper]{article}

% ============================================================
%                    КОДИРОВКА И ЯЗЫК
% ============================================================

\usepackage[T2A]{fontenc}
\usepackage[utf8]{inputenc}
\usepackage[russian,english]{babel}
\usepackage{cmap}

% ============================================================
%                    СТРАНИЦА И ТЕКСТ
% ============================================================

\usepackage[
    left=20mm,
    right=20mm,
    top=20mm,
    bottom=22mm,
    headheight=15pt
]{geometry}

\usepackage{indentfirst}
\usepackage{setspace}
\usepackage{microtype}
\usepackage{parskip}
\usepackage{multicol}
\usepackage{lipsum}

% ============================================================
%                         МАТЕМАТИКА
% ============================================================

\usepackage{amsmath}
\usepackage{amssymb}
\usepackage{amsfonts}
\usepackage{mathtools}
\usepackage{bm}
\usepackage{cancel}
\usepackage{esint}
\usepackage{nicefrac}

% ============================================================
%                  ТАБЛИЦЫ И СПИСКИ
% ============================================================

\usepackage{array}
\usepackage{booktabs}
\usepackage{longtable}
\usepackage{tabularx}
\usepackage{multirow}
\usepackage{makecell}
\usepackage{enumitem}

% ============================================================
%                  ГРАФИКА И ДИАГРАММЫ
% ============================================================

\usepackage{xcolor}
\usepackage{graphicx}
\usepackage{float}
\usepackage{wrapfig}
\usepackage{caption}
\usepackage{subcaption}

\usepackage{tikz}
\usetikzlibrary{
    arrows.meta,
    calc,
    positioning,
    angles,
    quotes,
    patterns,
    shapes.geometric,
    decorations.pathmorphing,
    decorations.markings
}

\usepackage{pgfplots}
\pgfplotsset{compat=1.18}

\usepackage{circuitikz}

% ============================================================
%                  ФИЗИКА И ЕДИНИЦЫ
% ============================================================

\usepackage{siunitx}
\sisetup{
    output-decimal-marker={,},
    per-mode=symbol,
    exponent-product=\cdot
}

% ============================================================
%                  РАМКИ И ОФОРМЛЕНИЕ
% ============================================================

\usepackage[most]{tcolorbox}

\definecolor{mainblue}{RGB}{42,91,160}
\definecolor{darkblue}{RGB}{24,48,85}
\definecolor{lightblue}{RGB}{229,240,255}
\definecolor{lightgreen}{RGB}{229,247,235}
\definecolor{lightred}{RGB}{255,232,232}
\definecolor{lightyellow}{RGB}{255,248,214}
\definecolor{codegray}{RGB}{245,245,245}
\definecolor{darkgray}{RGB}{70,70,70}
\definecolor{purple}{RGB}{115,66,170}

\newtcolorbox{definitionbox}[1][]{
    enhanced,
    breakable,
    colback=lightblue,
    colframe=mainblue,
    fonttitle=\bfseries,
    title={Определение},
    boxrule=0.8pt,
    arc=2mm,
    #1
}

\newtcolorbox{examplebox}[1][]{
    enhanced,
    breakable,
    colback=lightgreen,
    colframe=green!50!black,
    fonttitle=\bfseries,
    title={Пример},
    boxrule=0.8pt,
    arc=2mm,
    #1
}

\newtcolorbox{warningbox}[1][]{
    enhanced,
    breakable,
    colback=lightred,
    colframe=red!70!black,
    fonttitle=\bfseries,
    title={Важно},
    boxrule=0.8pt,
    arc=2mm,
    #1
}

\newtcolorbox{taskbox}[2][]{
    enhanced,
    breakable,
    colback=lightyellow,
    colframe=orange!70!black,
    fonttitle=\bfseries,
    title={Задача #2},
    boxrule=0.8pt,
    arc=2mm,
    #1
}

% ============================================================
%                         ТЕОРЕМЫ
% ============================================================

\usepackage{amsthm}

\newtheorem{theorem}{Теорема}[section]
\newtheorem{lemma}[theorem]{Лемма}
\newtheorem{corollary}[theorem]{Следствие}
\newtheorem{proposition}[theorem]{Утверждение}

\theoremstyle{definition}
\newtheorem{definition}[theorem]{Определение}
\newtheorem{example}[theorem]{Пример}

\theoremstyle{remark}
\newtheorem{remark}[theorem]{Замечание}

% ============================================================
%                     ПРОГРАММНЫЙ КОД
% ============================================================

\usepackage{listings}

\lstdefinestyle{customcode}{
    backgroundcolor=\color{codegray},
    basicstyle=\ttfamily\small,
    keywordstyle=\color{mainblue}\bfseries,
    commentstyle=\color{green!45!black},
    stringstyle=\color{purple},
    numberstyle=\tiny\color{darkgray},
    numbers=left,
    numbersep=8pt,
    frame=single,
    rulecolor=\color{black!20},
    breaklines=true,
    showstringspaces=false,
    tabsize=4,
    columns=fullflexible
}

\lstset{style=customcode}

% ============================================================
%                    ССЫЛКИ И ЗАГОЛОВКИ
% ============================================================

\usepackage{hyperref}

\hypersetup{
    colorlinks=true,
    linkcolor=mainblue,
    urlcolor=purple,
    citecolor=green!45!black,
    pdftitle={Мега-тестовый документ LaTeX},
    pdfauthor={ReNothingg},
    pdfsubject={Проверка возможностей LaTeX},
    pdfkeywords={LaTeX, TeX, математика, физика, TikZ, тест}
}

\usepackage{fancyhdr}
\usepackage{lastpage}

\pagestyle{fancy}
\fancyhf{}
\fancyhead[L]{\textbf{Мега-тест LaTeX}}
\fancyhead[R]{ReNothingg}
\fancyfoot[L]{Тестовый документ}
\fancyfoot[C]{Страница \thepage\ из \pageref{LastPage}}
\fancyfoot[R]{\today}

\renewcommand{\headrulewidth}{0.4pt}
\renewcommand{\footrulewidth}{0.4pt}

% ============================================================
%                    ПОЛЬЗОВАТЕЛЬСКИЕ КОМАНДЫ
% ============================================================

\newcommand{\R}{\mathbb{R}}
\newcommand{\N}{\mathbb{N}}
\newcommand{\Z}{\mathbb{Z}}
\newcommand{\Q}{\mathbb{Q}}
\newcommand{\C}{\mathbb{C}}

\newcommand{\vect}[1]{\bm{#1}}
\newcommand{\abs}[1]{\left|#1\right|}
\newcommand{\norm}[1]{\left\lVert#1\right\rVert}
\newcommand{\degree}{^{\circ}}
\newcommand{\dd}{\,\mathrm{d}}
\newcommand{\e}{\mathrm{e}}
\newcommand{\ii}{\mathrm{i}}

\DeclareMathOperator{\arctg}{arctg}
\DeclareMathOperator{\arcctg}{arcctg}
\DeclareMathOperator{\rank}{rank}
\DeclareMathOperator{\grad}{grad}
\DeclareMathOperator{\divergence}{div}
\DeclareMathOperator{\rot}{rot}

% ============================================================
%                         ДОКУМЕНТ
% ============================================================

\begin{document}

% ============================================================
%                      ТИТУЛЬНАЯ СТРАНИЦА
% ============================================================

\begin{titlepage}
    \centering

    \vspace*{1.5cm}

    {\Huge\bfseries МЕГА-ТЕСТОВЫЙ ДОКУМЕНТ\par}

    \vspace{0.4cm}

    {\LARGE для проверки возможностей \LaTeX\par}

    \vspace{1.5cm}

    \begin{tikzpicture}
        \draw[
            rounded corners=8pt,
            very thick,
            mainblue,
            fill=lightblue
        ]
        (0,0) rectangle (11,5);

        \node[
            font=\Huge\bfseries,
            text=darkblue
        ] at (5.5,3.2) {\LaTeX};

        \node[
            font=\Large,
            text=darkgray
        ] at (5.5,2.1) {математика, физика, таблицы, графики};

        \node[
            font=\large,
            text=darkgray
        ] at (5.5,1.3) {TikZ, PGFPlots, схемы, код и оформление};
    \end{tikzpicture}

    \vfill

    \begin{tabular}{rl}
        \textbf{Автор:} & ReNothingg \\[4pt]
        \textbf{Назначение:} & проверка рендеринга TeX \\[4pt]
        \textbf{Компилятор:} & pdfLaTeX \\[4pt]
        \textbf{Дата:} & \today
    \end{tabular}

    \vfill

    {\large Огромный демонстрационный документ\par}

\end{titlepage}

% ============================================================
%                         ОГЛАВЛЕНИЕ
% ============================================================

\tableofcontents
\newpage

\listoffigures
\listoftables
\newpage

% ============================================================
%                     ОСНОВНОЙ ТЕКСТ
% ============================================================

\section{Проверка обычного текста}

Это обычный русский текст, предназначенный для проверки кодировки,
переносов строк, абзацев, шрифтов и выравнивания. Документ содержит
кириллические символы: ё, й, ц, щ, ъ, ы, ь, э, ю, я.

Также проверяются кавычки: <<русские кавычки>>, ``английские кавычки'',
дефис -, короткое тире -- и длинное тире --- между частями предложения.

\textbf{Полужирный текст},
\textit{курсивный текст},
\underline{подчёркнутый текст},
\texttt{моноширинный текст},
\textsc{капитель},
\emph{логическое выделение}.

{\small Маленький текст.}

{\large Большой текст.}

{\Large Ещё более крупный текст.}

{\huge Очень крупный текст.}

\begin{center}
    Текст по центру.
\end{center}

\begin{flushright}
    Текст по правому краю.
\end{flushright}

\begin{flushleft}
    Текст по левому краю.
\end{flushleft}

\begin{quote}
Это пример цитаты. Цитата автоматически получает отступы слева и справа.
Она может содержать несколько предложений и переноситься на новые строки.
\end{quote}

\begin{quotation}
Это более длинный вариант цитирования. Окружение \texttt{quotation}
может использоваться для оформления фрагментов текста, состоящих из
нескольких абзацев.

Второй абзац внутри цитаты также получает специальное форматирование.
\end{quotation}

\section{Списки}

\subsection{Маркированный список}

\begin{itemize}
    \item Первый элемент.
    \item Второй элемент.
    \item Третий элемент:
    \begin{itemize}
        \item вложенный элемент;
        \item ещё один вложенный элемент;
        \begin{itemize}
            \item третий уровень вложенности.
        \end{itemize}
    \end{itemize}
\end{itemize}

\subsection{Нумерованный список}

\begin{enumerate}
    \item Подготовить исходный файл.
    \item Запустить компилятор.
    \item Исправить ошибки.
    \item Повторно скомпилировать документ.
\end{enumerate}

\subsection{Настраиваемый список}

\begin{enumerate}[label=\textbf{Шаг \arabic*:}, leftmargin=25mm]
    \item Установить дистрибутив TeX.
    \item Создать файл с расширением \texttt{.tex}.
    \item Добавить преамбулу.
    \item Начать окружение документа.
\end{enumerate}

\subsection{Список определений}

\begin{description}
    \item[\LaTeX] Система подготовки документов.
    \item[TeX] Система компьютерной вёрстки, созданная Дональдом Кнутом.
    \item[pdfLaTeX] Компилятор, создающий PDF-файл.
    \item[TikZ] Пакет для создания векторной графики.
\end{description}

% ============================================================
%                         МАТЕМАТИКА
% ============================================================

\section{Базовые математические формулы}

Встроенная формула выглядит так: $E = mc^2$.

Ещё одна встроенная формула:
$a^2+b^2=c^2$.

Отдельная формула:

\[
    x_{1,2}
    =
    \frac{-b\pm\sqrt{b^2-4ac}}{2a}.
\]

Формула с номером:

\begin{equation}
    \label{eq:einstein}
    E=mc^2.
\end{equation}

На формулу \eqref{eq:einstein} можно сослаться из текста.

\subsection{Степени, корни и дроби}

\[
    x^2,\qquad
    x^{100},\qquad
    2^{-5},\qquad
    \sqrt{x},\qquad
    \sqrt[3]{x},\qquad
    \sqrt[n]{a}.
\]

\[
    \frac{a}{b},
    \qquad
    \dfrac{x+1}{x-1},
    \qquad
    \nicefrac{1}{2},
    \qquad
    \cfrac{1}{
        1+\cfrac{1}{
            2+\cfrac{1}{3}
        }
    }.
\]

\subsection{Суммы, произведения и пределы}

\[
    \sum_{k=1}^{n} k
    =
    \frac{n(n+1)}{2}.
\]

\[
    \sum_{k=0}^{\infty} q^k
    =
    \frac{1}{1-q},
    \qquad
    \abs{q}<1.
\]

\[
    \prod_{k=1}^{n} k=n!.
\]

\[
    \lim_{x\to 0}\frac{\sin x}{x}=1.
\]

\[
    \lim_{n\to\infty}
    \left(1+\frac{1}{n}\right)^n
    =
    \e.
\]

\subsection{Тригонометрия}

\[
    \sin^2 x+\cos^2 x=1.
\]

\[
    \sin(\alpha+\beta)
    =
    \sin\alpha\cos\beta
    +
    \cos\alpha\sin\beta.
\]

\[
    \cos(\alpha+\beta)
    =
    \cos\alpha\cos\beta
    -
    \sin\alpha\sin\beta.
\]

\[
    \tg x=\frac{\sin x}{\cos x},
    \qquad
    \ctg x=\frac{\cos x}{\sin x}.
\]

\[
    \arcsin x,\qquad
    \arccos x,\qquad
    \arctg x,\qquad
    \arcctg x.
\]

\subsection{Логарифмы}

\[
    \log_a(xy)=\log_a x+\log_a y.
\]

\[
    \log_a\frac{x}{y}
    =
    \log_a x-\log_a y.
\]

\[
    \log_a x
    =
    \frac{\ln x}{\ln a}.
\]

\[
    a^{\log_a x}=x.
\]

% ============================================================
%                    СИСТЕМЫ И ВЫРАВНИВАНИЕ
% ============================================================

\section{Системы уравнений и выравнивание}

\subsection{Система уравнений}

\[
\begin{cases}
    2x+3y=7,\\
    5x-y=9.
\end{cases}
\]

\subsection{Система неравенств}

\[
\begin{cases}
    x^2-5x+6\geq 0,\\
    x+1<4,\\
    x\in\R.
\end{cases}
\]

\subsection{Кусочно заданная функция}

\[
f(x)=
\begin{cases}
    x^2, & x<0,\\
    0,   & x=0,\\
    \sqrt{x}, & x>0.
\end{cases}
\]

\subsection{Несколько строк вычислений}

\begin{align}
    (a+b)^2
    &=a^2+2ab+b^2,\\
    (a-b)^2
    &=a^2-2ab+b^2,\\
    (a+b)(a-b)
    &=a^2-b^2.
\end{align}

\begin{align*}
    x^3-1
    &=(x-1)(x^2+x+1),\\
    x^3+1
    &=(x+1)(x^2-x+1).
\end{align*}

\subsection{Длинное преобразование}

\begin{align*}
    \int x\e^x\dd x
    &=
    x\e^x-\int\e^x\dd x\\
    &=
    x\e^x-\e^x+C\\
    &=
    \e^x(x-1)+C.
\end{align*}

% ============================================================
%                          МАТРИЦЫ
% ============================================================

\section{Матрицы и определители}

\[
A=
\begin{pmatrix}
    1 & 2 & 3\\
    4 & 5 & 6\\
    7 & 8 & 9
\end{pmatrix}.
\]

\[
B=
\begin{bmatrix}
    a & b\\
    c & d
\end{bmatrix}.
\]

\[
C=
\begin{vmatrix}
    a & b\\
    c & d
\end{vmatrix}
=
ad-bc.
\]

\[
D=
\begin{Vmatrix}
    x_1 & x_2\\
    y_1 & y_2
\end{Vmatrix}.
\]

\[
I_4=
\begin{pmatrix}
    1&0&0&0\\
    0&1&0&0\\
    0&0&1&0\\
    0&0&0&1
\end{pmatrix}.
\]

\[
\begin{pmatrix}
    a&b\\
    c&d
\end{pmatrix}^{-1}
=
\frac{1}{ad-bc}
\begin{pmatrix}
    d&-b\\
    -c&a
\end{pmatrix}.
\]

\[
\rank A=2.
\]

% ============================================================
%                       АНАЛИЗ
% ============================================================

\section{Производные и интегралы}

\subsection{Производные}

\[
    \frac{\dd}{\dd x}x^n=nx^{n-1}.
\]

\[
    \frac{\dd}{\dd x}\sin x=\cos x.
\]

\[
    \frac{\dd}{\dd x}\cos x=-\sin x.
\]

\[
    \frac{\dd}{\dd x}\e^x=\e^x.
\]

\[
    \frac{\dd}{\dd x}\ln x=\frac{1}{x}.
\]

\[
    (fg)'=f'g+fg'.
\]

\[
    \left(\frac{f}{g}\right)'
    =
    \frac{f'g-fg'}{g^2}.
\]

\[
    \frac{\partial f}{\partial x},
    \qquad
    \frac{\partial^2 f}{\partial x^2},
    \qquad
    \frac{\partial^2 f}{\partial x\partial y}.
\]

\subsection{Неопределённые интегралы}

\[
    \int x^n\dd x
    =
    \frac{x^{n+1}}{n+1}+C,
    \qquad n\neq-1.
\]

\[
    \int\frac{1}{x}\dd x
    =
    \ln\abs{x}+C.
\]

\[
    \int\sin x\dd x
    =
    -\cos x+C.
\]

\[
    \int\cos x\dd x
    =
    \sin x+C.
\]

\subsection{Определённые интегралы}

\[
    \int_0^1 x^2\dd x
    =
    \left.\frac{x^3}{3}\right|_0^1
    =
    \frac{1}{3}.
\]

\[
    \int_{-\infty}^{+\infty}
    \e^{-x^2}\dd x
    =
    \sqrt{\pi}.
\]

\subsection{Кратные и контурные интегралы}

\[
    \iint\limits_D
    f(x,y)\dd x\dd y.
\]

\[
    \iiint\limits_V
    f(x,y,z)\dd x\dd y\dd z.
\]

\[
    \oint\limits_C
    \vect{F}\cdot\dd\vect{r}.
\]

\[
    \oiint\limits_S
    \vect{E}\cdot\dd\vect{S}.
\]

\subsection{Дифференциальное уравнение}

\[
    y''+4y'+4y=0.
\]

Характеристическое уравнение:

\[
    \lambda^2+4\lambda+4=0.
\]

\[
    (\lambda+2)^2=0.
\]

Общее решение:

\[
    y(x)=(C_1+C_2x)\e^{-2x}.
\]

% ============================================================
%                      КОМПЛЕКСНЫЕ ЧИСЛА
% ============================================================

\section{Комплексные числа}

\[
    z=a+b\ii.
\]

\[
    \ii^2=-1.
\]

\[
    \overline{z}=a-b\ii.
\]

\[
    \abs{z}=\sqrt{a^2+b^2}.
\]

\[
    z=r(\cos\varphi+\ii\sin\varphi).
\]

Формула Эйлера:

\[
    \e^{\ii\varphi}
    =
    \cos\varphi+\ii\sin\varphi.
\]

Следовательно,

\[
    \e^{\ii\pi}+1=0.
\]

Формула Муавра:

\[
    z^n
    =
    r^n
    \left(
        \cos n\varphi
        +
        \ii\sin n\varphi
    \right).
\]

% ============================================================
%                      ВЕКТОРЫ
% ============================================================

\section{Векторы}

\[
    \vect{a}
    =
    \begin{pmatrix}
        a_x\\
        a_y\\
        a_z
    \end{pmatrix}.
\]

\[
    \vect{a}+\vect{b}
    =
    \begin{pmatrix}
        a_x+b_x\\
        a_y+b_y\\
        a_z+b_z
    \end{pmatrix}.
\]

\[
    \vect{a}\cdot\vect{b}
    =
    a_xb_x+a_yb_y+a_zb_z.
\]

\[
    \vect{a}\times\vect{b}
    =
    \begin{vmatrix}
        \vect{i}&\vect{j}&\vect{k}\\
        a_x&a_y&a_z\\
        b_x&b_y&b_z
    \end{vmatrix}.
\]

\[
    \cos\alpha
    =
    \frac{\vect{a}\cdot\vect{b}}
    {\norm{\vect{a}}\norm{\vect{b}}}.
\]

\[
    \grad f
    =
    \nabla f.
\]

\[
    \divergence\vect{F}
    =
    \nabla\cdot\vect{F}.
\]

\[
    \rot\vect{F}
    =
    \nabla\times\vect{F}.
\]

% ============================================================
%                         ТЕОРЕМЫ
% ============================================================

\section{Теоремы, доказательства и определения}

\begin{definition}
Функция $f(x)$ называется непрерывной в точке $x_0$, если

\[
    \lim_{x\to x_0}f(x)=f(x_0).
\]
\end{definition}

\begin{theorem}[Теорема Пифагора]
В прямоугольном треугольнике квадрат гипотенузы равен сумме квадратов
катетов:

\[
    c^2=a^2+b^2.
\]
\end{theorem}

\begin{proof}
Построим квадрат со стороной $a+b$ и разместим внутри него четыре
равных прямоугольных треугольника. Сравнение площадей даёт

\[
    (a+b)^2=4\cdot\frac{ab}{2}+c^2.
\]

Раскрывая скобки, получаем

\[
    a^2+2ab+b^2=2ab+c^2,
\]

откуда

\[
    c^2=a^2+b^2.
\]
\end{proof}

\begin{lemma}
Если число делится на $4$, то оно является чётным.
\end{lemma}

\begin{proof}
Пусть $n=4k$, где $k\in\Z$. Тогда

\[
    n=2(2k),
\]

следовательно, $n$ делится на $2$.
\end{proof}

\begin{corollary}
Если квадрат натурального числа нечётный, то и само число нечётное.
\end{corollary}

\begin{remark}
Нумерация теорем, лемм и следствий связана с номером раздела.
\end{remark}

% ============================================================
%                         РАМКИ
% ============================================================

\section{Цветные рамки}

\begin{definitionbox}
Скорость --- это физическая величина, характеризующая быстроту изменения
положения тела в пространстве.

\[
    \vect{v}
    =
    \frac{\Delta\vect{r}}{\Delta t}.
\]
\end{definitionbox}

\begin{examplebox}
Автомобиль проехал расстояние $\SI{120}{\kilo\meter}$ за
$\SI{2}{\hour}$. Средняя скорость равна

\[
    v
    =
    \frac{s}{t}
    =
    \frac{\SI{120}{\kilo\meter}}{\SI{2}{\hour}}
    =
    \SI{60}{\kilo\meter\per\hour}.
\]
\end{examplebox}

\begin{warningbox}
При подстановке значений в физические формулы необходимо привести все
величины к одной системе единиц, обычно к системе СИ.
\end{warningbox}

\begin{taskbox}{1}
Тело начинает движение из состояния покоя с постоянным ускорением
$a=\SI{2}{\meter\per\second\squared}$. Найдите путь, пройденный за
$t=\SI{5}{\second}$.

\[
    s
    =
    v_0t+\frac{at^2}{2}
    =
    0+\frac{2\cdot5^2}{2}
    =
    \SI{25}{\meter}.
\]
\end{taskbox}

% ============================================================
%                         ФИЗИКА
% ============================================================

\section{Механика}

\subsection{Кинематика}

Равномерное движение:

\[
    x=x_0+vt.
\]

Равноускоренное движение:

\[
    v=v_0+at.
\]

\[
    x=x_0+v_0t+\frac{at^2}{2}.
\]

\[
    v^2-v_0^2=2as.
\]

Движение тела, брошенного под углом к горизонту:

\[
    x(t)=v_0\cos\alpha\cdot t.
\]

\[
    y(t)=v_0\sin\alpha\cdot t-\frac{gt^2}{2}.
\]

Максимальная высота:

\[
    h_{\max}
    =
    \frac{v_0^2\sin^2\alpha}{2g}.
\]

Дальность полёта:

\[
    L
    =
    \frac{v_0^2\sin2\alpha}{g}.
\]

\subsection{Динамика}

Второй закон Ньютона:

\[
    \sum\vect{F}=m\vect{a}.
\]

Сила тяжести:

\[
    \vect{F}_{\text{тяж}}=m\vect{g}.
\]

Сила упругости:

\[
    \vect{F}_{\text{упр}}=-k\vect{x}.
\]

Сила трения:

\[
    F_{\text{тр}}=\mu N.
\]

Закон всемирного тяготения:

\[
    F
    =
    G\frac{m_1m_2}{r^2}.
\]

\subsection{Работа и энергия}

Работа постоянной силы:

\[
    A=Fs\cos\alpha.
\]

Кинетическая энергия:

\[
    E_k=\frac{mv^2}{2}.
\]

Потенциальная энергия:

\[
    E_p=mgh.
\]

Энергия упругой деформации:

\[
    E_{\text{упр}}
    =
    \frac{kx^2}{2}.
\]

Закон сохранения механической энергии:

\[
    E_{k1}+E_{p1}
    =
    E_{k2}+E_{p2}.
\]

Мощность:

\[
    P=\frac{A}{t}.
\]

\subsection{Импульс}

\[
    \vect{p}=m\vect{v}.
\]

\[
    \vect{F}\Delta t=\Delta\vect{p}.
\]

Закон сохранения импульса:

\[
    m_1\vect{v}_1+m_2\vect{v}_2
    =
    m_1\vect{u}_1+m_2\vect{u}_2.
\]

\section{Молекулярная физика}

Количество вещества:

\[
    \nu=\frac{m}{M}.
\]

Количество частиц:

\[
    N=\nu N_A.
\]

Уравнение Менделеева--Клапейрона:

\[
    pV=\nu RT.
\]

Основное уравнение молекулярно-кинетической теории:

\[
    p
    =
    \frac{1}{3}m_0n\overline{v^2}.
\]

Средняя кинетическая энергия молекулы:

\[
    \overline{E_k}
    =
    \frac{3}{2}kT.
\]

Первый закон термодинамики:

\[
    Q=\Delta U+A.
\]

Коэффициент полезного действия:

\[
    \eta
    =
    \frac{A_{\text{полезная}}}
    {A_{\text{затраченная}}}
    \cdot100\%.
\]

\section{Электричество и магнетизм}

\subsection{Электростатика}

Закон Кулона:

\[
    F
    =
    k\frac{\abs{q_1q_2}}{r^2}.
\]

Напряжённость электрического поля:

\[
    \vect{E}
    =
    \frac{\vect{F}}{q}.
\]

Для точечного заряда:

\[
    E
    =
    k\frac{\abs{q}}{r^2}.
\]

Потенциал:

\[
    \varphi
    =
    \frac{W_p}{q}.
\]

Напряжение:

\[
    U=\varphi_1-\varphi_2.
\]

Электроёмкость:

\[
    C=\frac{q}{U}.
\]

Энергия конденсатора:

\[
    W
    =
    \frac{CU^2}{2}
    =
    \frac{qU}{2}
    =
    \frac{q^2}{2C}.
\]

\subsection{Постоянный ток}

Сила тока:

\[
    I=\frac{q}{t}.
\]

Закон Ома:

\[
    I=\frac{U}{R}.
\]

Сопротивление проводника:

\[
    R=\rho\frac{l}{S}.
\]

Последовательное соединение:

\[
    R=R_1+R_2+\dots+R_n.
\]

Параллельное соединение:

\[
    \frac{1}{R}
    =
    \frac{1}{R_1}
    +
    \frac{1}{R_2}
    +
    \dots
    +
    \frac{1}{R_n}.
\]

Работа электрического тока:

\[
    A=UIt.
\]

Мощность электрического тока:

\[
    P=UI=I^2R=\frac{U^2}{R}.
\]

Закон Джоуля--Ленца:

\[
    Q=I^2Rt.
\]

\subsection{Магнитное поле}

Сила Ампера:

\[
    F_A=BIl\sin\alpha.
\]

Сила Лоренца:

\[
    F_L=qvB\sin\alpha.
\]

Магнитный поток:

\[
    \Phi=B S\cos\alpha.
\]

Закон электромагнитной индукции:

\[
    \mathcal{E}_i
    =
    -\frac{\Delta\Phi}{\Delta t}.
\]

Энергия магнитного поля катушки:

\[
    W
    =
    \frac{LI^2}{2}.
\]

\section{Оптика}

Закон отражения:

\[
    \alpha=\beta.
\]

Закон преломления:

\[
    n_1\sin\alpha=n_2\sin\beta.
\]

Показатель преломления:

\[
    n=\frac{c}{v}.
\]

Формула тонкой линзы:

\[
    \frac{1}{F}
    =
    \frac{1}{d}
    +
    \frac{1}{f}.
\]

Оптическая сила:

\[
    D=\frac{1}{F}.
\]

Увеличение линзы:

\[
    \Gamma
    =
    \frac{H}{h}
    =
    \frac{f}{d}.
\]

Интерференционный максимум:

\[
    \Delta=m\lambda.
\]

Интерференционный минимум:

\[
    \Delta
    =
    \left(m+\frac{1}{2}\right)\lambda.
\]

% ============================================================
%                         ТАБЛИЦЫ
% ============================================================

\section{Таблицы}

\subsection{Простая таблица}

\begin{table}[H]
    \centering
    \caption{Основные физические величины}
    \label{tab:physics-values}

    \begin{tabular}{lll}
        \toprule
        \textbf{Величина} & \textbf{Обозначение} & \textbf{Единица СИ}\\
        \midrule
        Длина & $l$ & метр, $\si{\meter}$\\
        Масса & $m$ & килограмм, $\si{\kilogram}$\\
        Время & $t$ & секунда, $\si{\second}$\\
        Сила тока & $I$ & ампер, $\si{\ampere}$\\
        Температура & $T$ & кельвин, $\si{\kelvin}$\\
        \bottomrule
    \end{tabular}
\end{table}

Ссылка на таблицу: таблица \ref{tab:physics-values}.

\subsection{Таблица с вертикальными линиями}

\begin{table}[H]
    \centering
    \caption{Результаты тестирования}

    \begin{tabular}{|c|l|c|c|}
        \hline
        \textbf{№} &
        \textbf{Участник} &
        \textbf{Баллы} &
        \textbf{Оценка}\\
        \hline
        1 & Алексей & 95 & 10\\
        \hline
        2 & Мария & 87 & 9\\
        \hline
        3 & Павел & 100 & 10\\
        \hline
        4 & Анна & 76 & 8\\
        \hline
    \end{tabular}
\end{table}

\subsection{Объединённые ячейки}

\begin{table}[H]
    \centering
    \caption{Таблица с объединением ячеек}

    \begin{tabular}{|c|c|c|}
        \hline
        \multirow{2}{*}{Предмет} &
        \multicolumn{2}{c|}{Результат}\\
        \cline{2-3}
        & Балл & Оценка\\
        \hline
        Математика & 92 & 10\\
        \hline
        Физика & 88 & 9\\
        \hline
        Информатика & 100 & 10\\
        \hline
    \end{tabular}
\end{table}

\subsection{Таблица на всю ширину}

\begin{table}[H]
    \centering
    \caption{Таблица с автоматически изменяемой шириной}

    \begin{tabularx}{\textwidth}{|c|X|c|}
        \hline
        \textbf{№} &
        \textbf{Описание} &
        \textbf{Статус}\\
        \hline
        1 &
        Очень длинный текст, который автоматически переносится внутри
        ячейки таблицы и занимает доступную ширину страницы. &
        Готово\\
        \hline
        2 &
        Проверка работы окружения \texttt{tabularx} и столбца типа
        \texttt{X}. &
        Работает\\
        \hline
    \end{tabularx}
\end{table}

\subsection{Длинная таблица}

\begin{longtable}{|c|l|c|}
    \caption{Длинная таблица тестовых данных}\\
    \hline
    \textbf{№} & \textbf{Название} & \textbf{Значение}\\
    \hline
    \endfirsthead

    \hline
    \textbf{№} & \textbf{Название} & \textbf{Значение}\\
    \hline
    \endhead

    1 & Параметр A & 10\\ \hline
    2 & Параметр B & 20\\ \hline
    3 & Параметр C & 30\\ \hline
    4 & Параметр D & 40\\ \hline
    5 & Параметр E & 50\\ \hline
    6 & Параметр F & 60\\ \hline
    7 & Параметр G & 70\\ \hline
    8 & Параметр H & 80\\ \hline
    9 & Параметр I & 90\\ \hline
    10 & Параметр J & 100\\ \hline
    11 & Параметр K & 110\\ \hline
    12 & Параметр L & 120\\ \hline
    13 & Параметр M & 130\\ \hline
    14 & Параметр N & 140\\ \hline
    15 & Параметр O & 150\\ \hline
\end{longtable}

% ============================================================
%                         КОЛОНКИ
% ============================================================

\section{Текст в нескольких колонках}

\begin{multicols}{2}

\subsection*{Левая и правая колонки}

LaTeX автоматически распределяет текст между колонками. Такой режим
может использоваться для создания газет, шпаргалок, справочников,
учебных материалов и компактных тестовых документов.

Внутри колонок можно использовать формулы:

\[
    a^2+b^2=c^2.
\]

Также доступны списки:

\begin{itemize}
    \item первый пункт;
    \item второй пункт;
    \item третий пункт.
\end{itemize}

\columnbreak

Эта часть текста начинается после принудительного перехода в следующую
колонку с помощью команды \texttt{\textbackslash columnbreak}.

Например, формула площади круга:

\[
    S=\pi r^2.
\]

Длина окружности:

\[
    L=2\pi r.
\]

Объём шара:

\[
    V=\frac{4}{3}\pi r^3.
\]

\end{multicols}

% ============================================================
%                          TIKZ
% ============================================================

\section{Векторная графика TikZ}

\subsection{Простые геометрические фигуры}

\begin{figure}[H]
    \centering

    \begin{tikzpicture}[scale=1.1]
        \draw[very thick, mainblue] (0,0) rectangle (3,2);
        \draw[very thick, red] (5,1) circle (1);
        \draw[very thick, green!50!black]
            (8,0) -- (10,0) -- (9,2) -- cycle;

        \node at (1.5,-0.4) {Прямоугольник};
        \node at (5,-0.4) {Окружность};
        \node at (9,-0.4) {Треугольник};
    \end{tikzpicture}

    \caption{Простые фигуры, созданные с помощью TikZ}
\end{figure}

\subsection{Координатная система и вектор}

\begin{figure}[H]
    \centering

    \begin{tikzpicture}[scale=0.9]
        \draw[step=1cm, gray!20, very thin]
            (-4,-3) grid (5,4);

        \draw[-{Latex[length=3mm]}, thick]
            (-4,0) -- (5,0)
            node[right] {$x$};

        \draw[-{Latex[length=3mm]}, thick]
            (0,-3) -- (0,4)
            node[above] {$y$};

        \foreach \x in {-3,-2,-1,1,2,3,4}
            \draw (\x,0.08) -- (\x,-0.08)
            node[below] {\small $\x$};

        \foreach \y in {-2,-1,1,2,3}
            \draw (0.08,\y) -- (-0.08,\y)
            node[left] {\small $\y$};

        \draw[
            -{Latex[length=4mm]},
            very thick,
            mainblue
        ]
        (0,0) -- (3,2)
        node[midway, above left] {$\vect{a}$};

        \fill[mainblue] (3,2) circle (2pt);
        \node[above right] at (3,2) {$A(3;2)$};
    \end{tikzpicture}

    \caption{Координатная система и вектор}
\end{figure}

\subsection{Сложение двух сил}

\begin{figure}[H]
    \centering

    \begin{tikzpicture}[scale=1.1]
        \coordinate (O) at (0,0);
        \coordinate (A) at (4,0);
        \coordinate (B) at (2.2,3);
        \coordinate (C) at ($(A)+(B)$);

        \draw[
            -{Latex[length=4mm]},
            very thick,
            blue
        ]
        (O) -- (A)
        node[midway, below] {$\vect{F}_1$};

        \draw[
            -{Latex[length=4mm]},
            very thick,
            red
        ]
        (O) -- (B)
        node[midway, left] {$\vect{F}_2$};

        \draw[dashed, gray]
            (A) -- (C);

        \draw[dashed, gray]
            (B) -- (C);

        \draw[
            -{Latex[length=4mm]},
            ultra thick,
            green!50!black
        ]
        (O) -- (C)
        node[midway, above left] {$\vect{F}_{\text{рез}}$};

        \pic[
            draw,
            ->,
            "$\alpha$",
            angle eccentricity=1.4,
            angle radius=0.8cm
        ]
        {angle=A--O--B};
    \end{tikzpicture}

    \caption{Сложение векторов сил по правилу параллелограмма}
\end{figure}

Результирующая сила:

\[
    F_{\text{рез}}
    =
    \sqrt{
        F_1^2+
        F_2^2+
        2F_1F_2\cos\alpha
    }.
\]

\subsection{Блок-схема}

\begin{figure}[H]
    \centering

    \tikzset{
        startstop/.style={
            rectangle,
            rounded corners,
            minimum width=3cm,
            minimum height=1cm,
            text centered,
            draw=black,
            fill=red!20
        },
        process/.style={
            rectangle,
            minimum width=3.5cm,
            minimum height=1cm,
            text centered,
            draw=black,
            fill=blue!15
        },
        decision/.style={
            diamond,
            aspect=2,
            minimum width=3cm,
            minimum height=1cm,
            text centered,
            draw=black,
            fill=green!15
        },
        arrow/.style={
            thick,
            -{Latex[length=3mm]}
        }
    }

    \begin{tikzpicture}[node distance=1.5cm]
        \node (start) [startstop] {Начало};
        \node (input) [process, below=of start] {Ввести число $n$};
        \node (decision) [decision, below=of input] {$n>0$?};
        \node (positive) [process, below left=1.5cm and 1cm of decision]
            {Число положительное};
        \node (negative) [process, below right=1.5cm and 1cm of decision]
            {Число неположительное};
        \node (stop) [startstop, below=3.5cm of decision] {Конец};

        \draw [arrow] (start) -- (input);
        \draw [arrow] (input) -- (decision);

        \draw [arrow]
            (decision) -- node[above left] {да} (positive);

        \draw [arrow]
            (decision) -- node[above right] {нет} (negative);

        \draw [arrow] (positive) |- (stop);
        \draw [arrow] (negative) |- (stop);
    \end{tikzpicture}

    \caption{Пример блок-схемы}
\end{figure}

\subsection{Пружинный маятник}

\begin{figure}[H]
    \centering

    \begin{tikzpicture}[scale=1]
        \draw[very thick] (-2,3) -- (2,3);

        \foreach \x in {-1.8,-1.4,...,1.8}
            \draw[gray] (\x,3) -- ++(-0.25,0.25);

        \draw[
            thick,
            decoration={
                coil,
                aspect=0.45,
                segment length=4mm,
                amplitude=4mm
            },
            decorate
        ]
        (0,3) -- (0,0.8);

        \draw[
            thick,
            fill=mainblue!25
        ]
        (-0.8,0.8) rectangle (0.8,-0.2);

        \node at (0,0.3) {$m$};

        \draw[
            -{Latex[length=3mm]},
            thick,
            red
        ]
        (1.3,0.8) -- (1.3,-0.8)
        node[right] {$\vect{F}_{\text{тяж}}$};

        \draw[
            -{Latex[length=3mm]},
            thick,
            green!50!black
        ]
        (-1.3,0) -- (-1.3,1.6)
        node[left] {$\vect{F}_{\text{упр}}$};
    \end{tikzpicture}

    \caption{Пружинный маятник}
\end{figure}

% ============================================================
%                         PGFPLOTS
% ============================================================

\section{Графики функций}

\subsection{Парабола}

\begin{figure}[H]
    \centering

    \begin{tikzpicture}
        \begin{axis}[
            width=12cm,
            height=8cm,
            axis lines=middle,
            xlabel={$x$},
            ylabel={$y$},
            xmin=-4,
            xmax=4,
            ymin=-2,
            ymax=10,
            grid=both,
            minor tick num=1,
            samples=200,
            legend pos=north west
        ]
            \addplot[
                mainblue,
                very thick,
                domain=-3.2:3.2
            ]
            {x^2};

            \addlegendentry{$y=x^2$}
        \end{axis}
    \end{tikzpicture}

    \caption{График квадратичной функции}
\end{figure}

\subsection{Несколько функций}

\begin{figure}[H]
    \centering

    \begin{tikzpicture}
        \begin{axis}[
            width=13cm,
            height=8cm,
            axis lines=middle,
            xlabel={$x$},
            ylabel={$y$},
            xmin=-6.5,
            xmax=6.5,
            ymin=-1.5,
            ymax=1.5,
            grid=major,
            samples=300,
            legend style={
                at={(0.5,-0.18)},
                anchor=north,
                legend columns=3
            }
        ]
            \addplot[
                blue,
                very thick,
                domain=-6.28:6.28
            ]
            {sin(deg(x))};

            \addlegendentry{$y=\sin x$}

            \addplot[
                red,
                very thick,
                domain=-6.28:6.28
            ]
            {cos(deg(x))};

            \addlegendentry{$y=\cos x$}

            \addplot[
                green!50!black,
                thick,
                dashed,
                domain=-6.28:6.28
            ]
            {0.5*sin(deg(2*x))};

            \addlegendentry{$y=\frac12\sin2x$}
        \end{axis}
    \end{tikzpicture}

    \caption{Графики тригонометрических функций}
\end{figure}

\subsection{Экспонента и логарифм}

\begin{figure}[H]
    \centering

    \begin{tikzpicture}
        \begin{axis}[
            width=13cm,
            height=8cm,
            axis lines=middle,
            xlabel={$x$},
            ylabel={$y$},
            xmin=-3,
            xmax=5,
            ymin=-3,
            ymax=8,
            grid=both,
            samples=250,
            legend pos=north west
        ]
            \addplot[
                purple,
                very thick,
                domain=-3:2
            ]
            {exp(x)};

            \addlegendentry{$y=\e^x$}

            \addplot[
                orange,
                very thick,
                domain=0.05:5
            ]
            {ln(x)};

            \addlegendentry{$y=\ln x$}
        \end{axis}
    \end{tikzpicture}

    \caption{Экспоненциальная и логарифмическая функции}
\end{figure}

\subsection{Точечный график}

\begin{figure}[H]
    \centering

    \begin{tikzpicture}
        \begin{axis}[
            width=12cm,
            height=7cm,
            xlabel={Время, $\si{\second}$},
            ylabel={Скорость, $\si{\meter\per\second}$},
            xmin=0,
            xmax=7,
            ymin=0,
            ymax=15,
            grid=major
        ]
            \addplot[
                only marks,
                mark=*,
                mark size=3pt
            ]
            coordinates {
                (0,0)
                (1,2)
                (2,4.2)
                (3,6)
                (4,8.1)
                (5,10.2)
                (6,12)
            };

            \addplot[
                thick,
                dashed,
                domain=0:6
            ]
            {2*x};
        \end{axis}
    \end{tikzpicture}

    \caption{Экспериментальные точки и теоретическая зависимость}
\end{figure}

% ============================================================
%                      ЭЛЕКТРИЧЕСКАЯ СХЕМА
% ============================================================

\section{Электрическая схема}

\begin{figure}[H]
    \centering

    \begin{circuitikz}[american currents]
        \draw
        (0,0)
        to[battery1, l=$\mathcal{E}$]
        (0,4)
        to[switch, l=$K$]
        (3,4)
        to[R, l=$R_1$]
        (6,4)
        to[R, l=$R_2$]
        (6,0)
        -- (0,0);

        \draw
        (3,4)
        to[ammeter, l=$A$]
        (3,2)
        to[lamp, l=$L$]
        (6,2);

        \draw
        (6,4)
        to[voltmeter, l=$V$]
        (6,2);
    \end{circuitikz}

    \caption{Пример электрической цепи}
\end{figure}

Для последовательного участка:

\[
    R_{\text{общ}}=R_1+R_2.
\]

Сила тока:

\[
    I
    =
    \frac{\mathcal{E}}
    {R_{\text{общ}}+r}.
\]

% ============================================================
%                         ПРОГРАММНЫЙ КОД
% ============================================================

\section{Листинги программного кода}

\subsection{Python}

\begin{lstlisting}[language=Python, caption={Решение квадратного уравнения}]
import math


def solve_quadratic(a: float, b: float, c: float):
    if a == 0:
        raise ValueError("Coefficient a must not be zero")

    discriminant = b ** 2 - 4 * a * c

    if discriminant < 0:
        return []

    if discriminant == 0:
        return [-b / (2 * a)]

    sqrt_d = math.sqrt(discriminant)

    x1 = (-b + sqrt_d) / (2 * a)
    x2 = (-b - sqrt_d) / (2 * a)

    return [x1, x2]


print(solve_quadratic(1, -5, 6))
\end{lstlisting}

\subsection{C++}

\begin{lstlisting}[language=C++, caption={Пример программы на C++}]
#include <iostream>
#include <vector>
#include <algorithm>

int main() {
    std::vector<int> numbers = {5, 3, 8, 1, 2};

    std::sort(numbers.begin(), numbers.end());

    for (const int value : numbers) {
        std::cout << value << ' ';
    }

    std::cout << '\n';
    return 0;
}
\end{lstlisting}

\subsection{JavaScript}

\begin{lstlisting}[language=JavaScript, caption={Асинхронный запрос}]
async function loadData() {
    try {
        const response = await fetch("/api/data");

        if (!response.ok) {
            throw new Error(`HTTP error: ${response.status}`);
        }

        const data = await response.json();
        console.log(data);
    } catch (error) {
        console.error("Failed to load data:", error);
    }
}

loadData();
\end{lstlisting}

\subsection{HTML}

\begin{lstlisting}[language=HTML, caption={Простая HTML-страница}]
<!DOCTYPE html>
<html lang="ru">
<head>
    <meta charset="UTF-8">
    <meta name="viewport"
          content="width=device-width, initial-scale=1.0">
    <title>Тестовая страница</title>
</head>
<body>
    <main>
        <h1>Привет, мир!</h1>
        <p>Это тестовая HTML-страница.</p>
    </main>
</body>
</html>
\end{lstlisting}

% ============================================================
%                  ЗАДАЧИ С РЕШЕНИЯМИ
% ============================================================

\section{Большой набор тестовых задач}

\subsection{Алгебра}

\begin{taskbox}{2}
Решите уравнение:

\[
    x^2-5x+6=0.
\]

Разложим выражение на множители:

\[
    x^2-5x+6
    =
    (x-2)(x-3).
\]

Следовательно,

\[
    x_1=2,
    \qquad
    x_2=3.
\]
\end{taskbox}

\begin{taskbox}{3}
Упростите выражение:

\[
    \frac{x^2-9}{x^2-6x+9}.
\]

Разложим числитель и знаменатель:

\[
    x^2-9=(x-3)(x+3),
\]

\[
    x^2-6x+9=(x-3)^2.
\]

Тогда при $x\neq3$:

\[
    \frac{x^2-9}{x^2-6x+9}
    =
    \frac{x+3}{x-3}.
\]
\end{taskbox}

\begin{taskbox}{4}
Решите неравенство:

\[
    \frac{x-2}{x+1}\geq0.
\]

Критические точки:

\[
    x=-1,
    \qquad
    x=2.
\]

Ответ:

\[
    x\in(-\infty,-1)\cup[2,+\infty).
\]
\end{taskbox}

\subsection{Геометрия}

\begin{taskbox}{5}
В прямоугольном треугольнике катеты равны
$a=\SI{6}{\centi\meter}$ и $b=\SI{8}{\centi\meter}$.
Найдите гипотенузу.

\[
    c
    =
    \sqrt{a^2+b^2}
    =
    \sqrt{6^2+8^2}
    =
    \sqrt{100}
    =
    \SI{10}{\centi\meter}.
\]
\end{taskbox}

\begin{taskbox}{6}
Найдите площадь треугольника со сторонами
$a=5$, $b=6$, $c=7$.

Полупериметр:

\[
    p
    =
    \frac{a+b+c}{2}
    =
    9.
\]

По формуле Герона:

\[
    S
    =
    \sqrt{p(p-a)(p-b)(p-c)}.
\]

\[
    S
    =
    \sqrt{9\cdot4\cdot3\cdot2}
    =
    \sqrt{216}
    =
    6\sqrt{6}.
\]
\end{taskbox}

\subsection{Физика}

\begin{taskbox}{7}
Тело массой $m=\SI{5}{\kilogram}$ движется с ускорением
$a=\SI{3}{\meter\per\second\squared}$. Найдите результирующую силу.

\[
    F=ma.
\]

\[
    F
    =
    \SI{5}{\kilogram}
    \cdot
    \SI{3}{\meter\per\second\squared}
    =
    \SI{15}{\newton}.
\]
\end{taskbox}

\begin{taskbox}{8}
К резистору сопротивлением $R=\SI{20}{\ohm}$ приложено напряжение
$U=\SI{12}{\volt}$. Найдите силу тока.

\[
    I
    =
    \frac{U}{R}
    =
    \frac{\SI{12}{\volt}}{\SI{20}{\ohm}}
    =
    \SI{0.6}{\ampere}.
\]
\end{taskbox}

\begin{taskbox}{9}
Тело брошено вертикально вверх со скоростью
$v_0=\SI{20}{\meter\per\second}$. Найдите максимальную высоту,
принимая $g=\SI{10}{\meter\per\second\squared}$.

В верхней точке $v=0$:

\[
    v^2=v_0^2-2gh.
\]

\[
    h
    =
    \frac{v_0^2}{2g}
    =
    \frac{20^2}{2\cdot10}
    =
    \SI{20}{\meter}.
\]
\end{taskbox}

% ============================================================
%                     СПЕЦИАЛЬНЫЕ СИМВОЛЫ
% ============================================================

\section{Специальные математические символы}

\[
\begin{gathered}
    \alpha,\beta,\gamma,\delta,\varepsilon,\zeta,\eta,\theta,\\
    \iota,\kappa,\lambda,\mu,\nu,\xi,\pi,\rho,\sigma,\tau,\\
    \upsilon,\varphi,\chi,\psi,\omega.
\end{gathered}
\]

\[
\begin{gathered}
    \Gamma,\Delta,\Theta,\Lambda,\Xi,\Pi,\Sigma,\Phi,\Psi,\Omega.
\end{gathered}
\]

\[
    \leq,\geq,\neq,\approx,\sim,\equiv,\propto,\ll,\gg.
\]

\[
    \in,\notin,\subset,\subseteq,\supset,\supseteq,\varnothing.
\]

\[
    \cup,\cap,\setminus,\times,\otimes,\oplus.
\]

\[
    \forall,\exists,\nexists,\Rightarrow,\Leftarrow,\Leftrightarrow.
\]

\[
    \infty,\partial,\nabla,\ell,\hbar,\Re,\Im.
\]

\[
    \leftarrow,\rightarrow,\leftrightarrow,
    \Leftarrow,\Rightarrow,\Leftrightarrow.
\]

\[
    \uparrow,\downarrow,\nearrow,\searrow,\swarrow,\nwarrow.
\]

% ============================================================
%                     СКОБКИ И РАЗМЕРЫ
% ============================================================

\section{Скобки различных типов}

\[
    (x),
    \quad
    [x],
    \quad
    \{x\},
    \quad
    \abs{x},
    \quad
    \norm{x}.
\]

\[
    \left(
        \frac{x^2+1}{x^2-1}
    \right).
\]

\[
    \left[
        \sum_{k=1}^{n}k^2
    \right].
\]

\[
    \left\{
        x\in\R
        \mid
        x^2<4
    \right\}.
\]

\[
    \left\langle
        \vect{a},\vect{b}
    \right\rangle.
\]

\[
    \left\lfloor x\right\rfloor,
    \qquad
    \left\lceil x\right\rceil.
\]

% ============================================================
%                      ЗАЧЁРКИВАНИЯ
% ============================================================

\section{Сокращения и зачёркивания}

\[
    \frac{\cancel{x}(x+1)}
    {\cancel{x}(x-1)}
    =
    \frac{x+1}{x-1},
    \qquad x\neq0.
\]

\[
    \cancelto{0}{x-x}+5=5.
\]

\[
    \frac{
        \bcancel{a}b
    }{
        \bcancel{a}c
    }
    =
    \frac{b}{c}.
\]

% ============================================================
%                       ССЫЛКИ
% ============================================================

\section{Гиперссылки и сноски}

Официальный сайт проекта можно указать как обычную ссылку:

\[
    \text{\url{https://example.com}}
\]

Можно сделать ссылкой произвольный текст:
\href{https://www.latex-project.org/}{официальный сайт проекта LaTeX}.

Это предложение содержит сноску.\footnote{
Сноски автоматически нумеруются и размещаются внизу страницы.
}

Ещё одна сноска.\footnote{
Вторая тестовая сноска с дополнительным текстом.
}

% ============================================================
%                     ТЕСТ ПЕРЕНОСОВ
% ============================================================

\section{Большой текстовый блок}

LaTeX является системой компьютерной вёрстки, предназначенной для
создания технических, научных, математических и учебных документов.
В отличие от обычных текстовых редакторов, пользователь описывает
структуру документа при помощи команд, после чего специальный
компилятор формирует итоговый PDF-файл.

Особенно хорошо LaTeX подходит для документов, содержащих большое
количество математических выражений. Система автоматически управляет
размерами символов, расстояниями между элементами формулы, нумерацией
уравнений, ссылками, сносками, оглавлением и библиографией.

При подготовке большого документа рекомендуется разделять исходный код
на логические части, использовать пользовательские команды для часто
повторяющихся конструкций и не смешивать содержание документа с его
визуальным оформлением. Большинство параметров оформления размещаются
в преамбуле, то есть до команды начала документа.

Документ может содержать главы, разделы, подразделы, таблицы, рисунки,
графики, формулы, программный код, цитаты, списки, теоремы,
доказательства и перекрёстные ссылки. Благодаря пакетам TikZ и PGFPlots
можно создавать сложную векторную графику непосредственно внутри
исходного файла.

% ============================================================
%                     КОНТРОЛЬНАЯ РАБОТА
% ============================================================

\section{Пример контрольной работы}

\begin{center}
    {\Large\bfseries Контрольная работа по математике и физике}

    \vspace{3mm}

    Вариант №1
\end{center}

\textbf{Фамилия и имя:}
\underline{\hspace{8cm}}

\vspace{4mm}

\textbf{Дата:}
\underline{\hspace{4cm}}

\vspace{8mm}

\subsection*{Часть A. Краткий ответ}

\begin{enumerate}[label=\textbf{A\arabic*.}]
    \item Решите уравнение:

    \[
        3x-7=11.
    \]

    \vspace{10mm}

    \item Вычислите:

    \[
        \sqrt{49}+\sqrt[3]{27}.
    \]

    \vspace{10mm}

    \item Упростите:

    \[
        \frac{x^2-4}{x-2}.
    \]

    \vspace{10mm}

    \item Найдите производную:

    \[
        f(x)=3x^4-2x^2+7.
    \]

    \vspace{10mm}
\end{enumerate}

\subsection*{Часть B. Полное решение}

\begin{enumerate}[label=\textbf{B\arabic*.}]
    \item Решите систему:

    \[
    \begin{cases}
        x+y=7,\\
        2x-y=5.
    \end{cases}
    \]

    \vspace{25mm}

    \item Найдите площадь фигуры, ограниченной графиками

    \[
        y=x
        \qquad\text{и}\qquad
        y=x^2.
    \]

    \vspace{30mm}

    \item Автомобиль начал движение из состояния покоя и за
    $\SI{10}{\second}$ достиг скорости
    $\SI{20}{\meter\per\second}$.

    Найдите ускорение и пройденный путь.

    \vspace{30mm}
\end{enumerate}

\subsection*{Часть C. Повышенный уровень}

\begin{enumerate}[label=\textbf{C\arabic*.}]
    \item Докажите, что для всех $n\in\N$ выполняется:

    \[
        1+2+\dots+n
        =
        \frac{n(n+1)}{2}.
    \]

    \vspace{35mm}

    \item Исследуйте функцию

    \[
        f(x)
        =
        \frac{x^2+1}{x-1}
    \]

    на область определения, точки экстремума, промежутки монотонности
    и асимптоты.

    \vspace{40mm}
\end{enumerate}

% ============================================================
%                    ТАБЛИЦА ОТВЕТОВ
% ============================================================

\section{Таблица ответов}

\begin{table}[H]
    \centering
    \renewcommand{\arraystretch}{2}

    \begin{tabularx}{\textwidth}{|c|X|c|}
        \hline
        \textbf{Задание} &
        \textbf{Ответ} &
        \textbf{Баллы}\\
        \hline
        A1 & & 1\\
        \hline
        A2 & & 1\\
        \hline
        A3 & & 1\\
        \hline
        A4 & & 1\\
        \hline
        B1 & & 3\\
        \hline
        B2 & & 4\\
        \hline
        B3 & & 4\\
        \hline
        C1 & & 5\\
        \hline
        C2 & & 5\\
        \hline
        \multicolumn{2}{|r|}{\textbf{Итого:}} & 25\\
        \hline
    \end{tabularx}
\end{table}

% ============================================================
%                         БИБЛИОГРАФИЯ
% ============================================================

\section{Пример библиографии}

В тексте можно сослаться на книгу Кнута \cite{knuth1984}
или руководство по LaTeX \cite{lamport1994}.

\begin{thebibliography}{9}

\bibitem{knuth1984}
Donald E. Knuth.
\textit{The TeXbook}.
Addison-Wesley, 1984.

\bibitem{lamport1994}
Leslie Lamport.
\textit{\LaTeX: A Document Preparation System}.
Addison-Wesley, 1994.

\bibitem{mittelbach2004}
Frank Mittelbach and Michel Goossens.
\textit{The \LaTeX\ Companion}.
Addison-Wesley, 2004.

\end{thebibliography}

% ============================================================
%                         ФИНАЛ
% ============================================================

\newpage

\begin{center}

    \vspace*{4cm}

    {\Huge\bfseries Документ успешно завершён}

    \vspace{1cm}

    \begin{tikzpicture}
        \draw[
            very thick,
            green!50!black,
            fill=lightgreen
        ]
        (0,0) circle (1.5);

        \draw[
            line width=3pt,
            green!50!black,
            rounded corners
        ]
        (-0.7,0) -- (-0.15,-0.55) -- (0.8,0.65);
    \end{tikzpicture}

    \vspace{1cm}

    {\Large Если эта страница отображается, основной тест пройден.}

    \vspace{1cm}

    \rule{0.7\textwidth}{0.5pt}

    \vspace{0.5cm}

    \textbf{ReNothingg --- Mega \LaTeX\ Test}

\end{center}

\end{document}